%Template da SBC para artigos em português

\documentclass[12pt]{article}

\usepackage{sbc-template}
\usepackage{graphicx,url}
\usepackage[utf8]{inputenc}
\usepackage[brazil]{babel}
%\usepackage[latin1]{inputenc}  

     
\sloppy

\title{Digital Sovereignty and the Economic Impact of DNS Centralization}

\author{Felipe S. Bertoglio\inst{1}}

\address{Instituto de Informática -- Universidade Federal do Rio Grande do Sul (UFRGS)\\
  Caixa Postal 15.064 -- 91.501-970 -- Porto Alegre -- RS -- Brazil\\
  \email{fsbertoglio@inf.ufrgs.br}
}

\begin{document} 

\maketitle

\begin{abstract}
  This paper examines digital sovereignty from an economic perspective, focusing on the centralization of Domain Name System (DNS) services among a few global providers. While existing literature addresses DNS centralization, the economic impacts on national sovereignty have not been explored. This study aims to fill that gap by creating a visualization panel to illustrate the economic distribution of DNS services in Brazilian companies and assessing the financial consequences of DNS failures on Brazil’s economy. Key issues include the economic impact of DNS service disruptions and the level of financial dependence of Brazilian firms on foreign providers.
\end{abstract}
     
\begin{resumo} 
  Este artigo examina a soberania digital a partir de uma perspectiva econômica, com foco na centralização dos serviços de Sistema de Nomes de Domínio (DNS) entre poucos provedores globais. Embora a literatura existente aborde a centralização do DNS, os impactos econômicos sobre a soberania nacional ainda não foram explorados. Este estudo busca preencher essa lacuna criando um painel de visualização para ilustrar a distribuição econômica dos serviços de DNS em empresas brasileiras e avaliando as consequências financeiras das falhas no DNS para a economia do Brasil. As questões-chave incluem o impacto econômico das interrupções nos serviços de DNS e o nível de dependência financeira das empresas brasileiras em relação a provedores estrangeiros.
\end{resumo}


\section{Introduction}

Sovereignty, particularly in the realm of computer networks, has become increasingly important in both academic discussions and the formulation of public policies, as well as in the positioning of individual countries. The study of digital sovereignty, as a contemporary interpretation of an age-old concept, examines various dimensions of computer networks within a continually evolving international political landscape. Key areas of focus in this field include infrastructure, device manufacturing, intellectual property, standards and regulations, security, and data flow. 

In this paper, we aim to analyze digital sovereignty from an economic perspective, concentrating on the centralization of Domain Name System (DNS) services under the control of a limited number of actors. This phenomenon has been extensively explored in the literature, where prior studies have introduced metrics and mechanisms to evaluate the market share held by various organizations, including thorough analyses of the national contexts in several countries. 

However, existing research lacks a specific focus on quantifying a nation's digital sovereignty through the economic and financial implications of DNS centralization. Thus, this work seeks to develop a visualization panel that depicts the economic distribution of DNS services among Brazilian companies and to assess the potential financial impact of DNS system failures on the Brazilian economy. 

Consequently, this study is driven by two primary motivations. First, considering the critical role of DNS services in the Brazilian economy, we examine the potential repercussions of service failures—whether due to malicious intent or accidental occurrences—on the economy. Second, we propose an analysis of the financial dependence of Brazilian companies on services provided by foreign entities, aiming to measure the cost of these services in relation to the revenue generated by the respective companies.



\section{Literature Review}

This chapter presents the concepts necessary to understand the work, as well as positioning it in the context of existing works that address digital sovereignty and DNS concentration. First, we list the publications that address the topic, presenting their contributions to the work. Then, we present the concept of sovereignty in the context of computer networks. We also explain what Domain Name Service is and how its concentration in the hands of a few organizations has been noted.

\subsection{Related Work}

To be done.

\subsection{Digital Sovereignty}

The concept of sovereignty in the context of computer networks is typically used to assert some form of collective control over digital content and infrastructures. However, interpretations, subjects, means, and definitions of sovereignty can vary significantly, as contextualized by \cite{couture:19}. The modern idea of sovereignty, articulated by thinkers like Machiavelli, Bodin, and Hobbes, refers to the supreme authority exercised over a territory or political entity. Until the 20th century, this concept was predominantly applied to control and authority within clearly defined territorial limits..
In today’s digital age, the notion of sovereignty has taken on new importance in public policy formulation and international relations. In the digital realm, the infosphere does not represent a physical territory, data is not finite, and digital assets primarily belong to private companies influenced by market forces and national governments. Therefore, the digital age compels us to reconsider the nature of sovereignty \cite{floridi:20}.

Various terms related to sovereignty have been analyzed in the digital sphere, including digital sovereignty, cyber sovereignty, internet sovereignty, technological sovereignty, and data sovereignty. \cite{azevedo:24} dopts the term "digital sovereignty" as it serves as an umbrella term in the current literature,covering multiple aspects.

As presented by \cite{floridi:20}, in the realm of computer networks, digital sovereignty refers to the control exercised over the physical layers (resources, infrastructure, devices), code (standards, rules, design), and information (content, data) of digital technologies. This control encompasses the ability to influence and restrict the production, design, use, and outcomes of these technologies.

Countries around the world are making significant efforts to regulate data privacy, enhance cybersecurity, and assert their national sovereignty in the technological sphere. As \cite{aydn:19} unlike the traditional conception of sovereignty, digital sovereignty lacks a consolidated international treaty and is approached differently by each nation based on its political and economic priorities.

The urgency of asserting digital sovereignty increased following Edward Snowden's revelations regarding extensive espionage and surveillance conducted by the United States. In the wake of these revelations, political leaders have utilized the concept of digital sovereignty to advocate for reducing technological dependence on foreign nations. A notable case is Brazil: after the Snowden allegations, then-President Dilma Rousseff proposed a plan to diminish the influence of the U.S. and its tech giants on the Brazilian internet. Similarly, Canada has invested in enhancing its digital infrastructure to decrease data traffic passing through U.S. networks \cite{couture:19}. The growing rivalry between the U.S. and China, along with the COVID-19 pandemic, has further accelerated global efforts to reduce technological and economic dependencies.

The European Union has been actively working to safeguard its digital sovereignty and to achieve strategic technological autonomy \cite{barrinha:22}. It addresses digital sovereignty through both technical and regulatory measures. On the technical side, this includes ensuring data confidentiality and integrity via standards-compliant hardware, the adoption of open-source software, and localized data processing and storage through the European Federated Cloud Initiative. On the regulatory side, it employs the General Data Protection Regulation (GDPR) \cite{tipper:22}.


\subsection{Domain Name System}
Virtually all modern internet applications depend on the Domain Name System (DNS) to function properly. Its main function, translating names into IP addresses, is essential, especially for users, who are accustomed to recognizing “locations” on the internet through friendly words, titles, and abbreviations, rather than numeric IP addresses \cite{deccio:09}.

Although the concept of name resolution is relatively straightforward, the system that implements it is highly complex and has broad impacts. Often, the name resolution of a domain depends on servers that are outside the control of the domain owner and are managed by third parties \cite{deccio:09}. The large scale of the database and the frequency of its updates make it essential to maintain this system in a distributed manner \cite{mockapetris:87a}.

According to \cite{tanenbaum:21}, the management of the top of the Internet's naming hierarchy is carried out by the Internet Corporation for Assigned Names and Numbers (ICANN). The conceptual structure of the Internet is organized into top-level domains (Top-Level Domains - TLDs), each housing several hosts. These TLDs are subdivided into subdomains, which can be continuously divided, forming a tree-like naming hierarchy, whose leaves represent domains that do not have subdomains.

A Top-Level Domain (TLD) is a zone located immediately below the root of the DNS hierarchy, exemplified by extensions such as ".com" or ".jp". TLDs are often classified into subgroups, such as Country Code Top-Level Domains (ccTLDs), which represent specific countries or regions (e.g., ".br" for Brazil), and Generic Top-Level Domains (gTLDs), which encompass generic categories such as ".org" or ".edu" \cite{hoffman:19}.

TLDs are operated by companies called registries, designated by ICANN. For example, Verisign is responsible for registering the ".com" domain. At a lower hierarchical level, registrars are responsible for selling domain names directly to end users. The market has several registrars, which compete with each other by offering different prices and services. Popular examples are: Domain.com, GoDaddy and NameCheap \cite{tanenbaum:21}. In Brazil, Registro.br is responsible for the registration and maintenance of domain names that use the .br \cite{registrobr}.

The DNS namespace is organized hierarchically. Zones represent parts of this namespace, are managed by a single entity, and are delegated to organizations in a descending manner, starting with the root zone. In the resolution process, name servers (resolvers) send queries to authoritative servers, which can provide a direct answer or a referral to a delegated zone \cite{deccio:09}. Based on data obtained from OpenIntel, between 2017 and 2021, \cite{zembruzki:22} quantified the numbers of TDLs, domains, NSs and ASs, for both IPv4 and IPv6.

As presented by \cite{azevedo:24}, an Autonomous System (AS) is a network of interconnected computing devices that operate under the same policy. It is typically managed by a single entity, such as an Internet Service Provider (ISP) or a technology organization, and is identified by an Autonomous System Number (ASN). Therefore, it is possible to associate the IP address of any Name Server (NS) with an AS, and consequently, with its operator and geographical region.

\subsection{DNS Centralization}
According to \cite{moura:20}, in recent years, there has been a growing concern about the centralization of the Internet, which consists of the concentration of services, infrastructure and users in a few companies. \cite{wang:21} when analyzing the 10 thousand most accessed sites on the tranco list \cite{lepochat:19}, concludes that the companies Cloudflare and Amazon alone concentrate more than 30\% of DNS domains.

When analyzing data from the \cite{openintel} project, and adding it to the prefix-to-AS data and AS-to-organization from CAIDA \cite{caida:13} \cite{caida:14}, \cite{zembruzki:22} obtained an overview of the evolution of DNS centralization over the period from 2017 to 2021. The study concludes that in 2021 the concentration of domains using IPv4 from the top 5 DNS ASs reached the 40\% mark, with 80\% of the domains concentrated in the top 100 ASes of DNS providers. 

Such concentration generates vulnerability, as presented by \cite{moura:20}, as it creates points of dependence susceptible to attacks, such as the Denial-of-Service (DoS) attacks against DNS service providers that occurred in 2016, harming Dyn \cite{hilton:16}, \cite{chacos:16}, and in 2019, harming AWS \cite{williams:19}. In both cases, important websites and services experienced outages.

A DNS outage occurs when a domain name cannot be resolved to its corresponding IP address. During this time, clients will send a DNS query for a domain name, but the DNS resolver may respond with an outdated IP address from its cache, which will not be functional. Alternatively, if the resolver attempts to query the authoritative name server for the domain and does not receive a response, the resolution will fail \cite{pramatarov:24}. 
A DNS outage occurs when a domain name cannot be resolved to its corresponding IP address. During this time, clients will send a DNS query for a domain name, but the DNS resolver may respond with an outdated IP address from its cache, which will not be functional. Alternatively, if the resolver attempts to query the authoritative name server for the domain and does not receive a response, the resolution will fail \cite{pramatarov:24}. The main causes of DNS failures include Distributed Denial of Service (DDoS) attacks, maintenance of authoritative name servers, issues in the data center hosting the authoritative name server, delays in DNS propagation, or misconfigurations.


\section{Proposed Work}

This chapter presents the methodology used to develop the work, detailing the steps taken to achieve the objectives. The methodology is divided into two main parts: the data collection and the development of the visualization panel and financial impact analysis.

\subsection{Data Collection}

The visualization panel will be developed using as inspiration the pipeline proposed by \cite{azevedo:24}, which determines the entire flow from the domain name to the organization handling the AS that manages the IP of the associated NS. The data will be collected from the OpenIntel project, which provides information on the DNS hierarchy, and from the CAIDA project, which provides information on the AS. Additionally, financial information will be added to the pipeline, such as the revenue generated by Brazilian companies that rely on DNS services provided by foreign entities, as well as the cost of these services. The implementation will be done using Python scripts, which will extract the necessary information and generate the visualization panel.
 

\subsection{Financial Impact Analysis}

The financial impact analysis will be conducted using the revenue generated by Brazilian companies that rely on DNS services provided by foreign entities, as well as the cost of these services. The model will simulate the financial consequences of DNS failures on the Brazilian economy, taking into account the number of companies affected, the duration of the outage, and the revenue lost due to the outage. 

\section{Schedule}

The work will be carried out over the next six months, or 24 weeks, and is divided into the following tasks:

\begin{itemize}
    \item Literature review and definition of the range of scope, weeks 1 to 4.
    \item Development of the pipeline for data collection, weeks 5 to 16.
    \item Development of the visualization painel and analysis of the results, weeks 10 to 18.
    \item Writing of the thesis, weeks 5 to 20.
    \item Final adjustments and submission of the work, weeks 21 to 24. 
    .
\end{itemize}

\section{Conclusion}

This report aims to analyze digital sovereignty from an economic perspective, focusing on the centralization of DNS services among a few global providers. While existing literature addresses DNS centralization, the economic impacts on national sovereignty have not been explored. This study aims to fill that gap by creating a visualization panel to illustrate the economic distribution of DNS services in Brazilian companies and assessing the financial consequences of DNS failures on Brazil’s economy. Key issues include the economic impact of DNS service disruptions and the level of financial dependence of Brazilian firms on foreign providers.

\bibliographystyle{sbc}
\bibliography{bibliography}

\end{document}
